\section{DECLARAÇÃO SOBRE USO DE INTELIGÊNCIA ARTIFICIAL }

Em conformidade com as diretrizes de integridade acadêmica e transparência na pesquisa científica, o autor declara o uso de ferramentas de Inteligência Artificial Generativa — especificamente o modelo de linguagem Google Gemini — como tecnologia assistiva durante as etapas de ideação, redação, formatação e revisão deste Trabalho de Conclusão de Curso.

O emprego da referida ferramenta ocorreu sob supervisão e direcionamento estrito do autor, restringindo-se às seguintes frentes operacionais:
\begin{itemize}
    \item \textbf{Geração Direta de Texto:} A ferramenta foi empregada para a consolidação textual do Resumo (e sua respectiva tradução para o \textit{Abstract}) e para a redação desta própria seção de Declaração de Uso de Inteligência Artificial. Destaca-se que a geração ocorreu estritamente a partir das premissas, dados analíticos e direcionamentos estruturais fornecidos previamente pelo autor.

    \item \textbf{Consultoria Metodológica, Normativa e Ética:} Realização de consultas referentes à conformidade do documento com as regras da Associação Brasileira de Normas Técnicas (ABNT), abrangendo o refatoramento de citações e a obtenção de guias orientativos sobre o escopo de conteúdo exigido para determinadas seções do trabalho.

    \item \textbf{Revisão Gramatical e Refinamento Estilístico:} Lapidação do texto em língua portuguesa, auxiliando na adequação da sintaxe, eliminação de vícios de linguagem, garantia de coesão textual e elevação do vocabulário para manter o tom formal, objetivo e impessoal exigido pela escrita acadêmica.

    \item \textbf{Codificação LaTeX e Visualização de Dados:} A ferramenta atuou como suporte técnico na formatação estrutural do documento e na resolução de problemas de compilação, assegurando a conformidade com as normas da ABNT. O auxílio abrangeu a organização modular do código-fonte e a criação do \textit{boxplot} utilizado na apresentação dos dados na seção \ref{sec:numeros_hell_yeah}.

    \item \textbf{Discussão Conceitual e Ideação:} Uso do modelo de linguagem como ferramenta de \textit{brainstorming}, englobando a formulação de ideias de títulos para o trabalho e a organização em tópicos de determinadas seções --- a exemplo da separação encontrada na seção de resultados.
\end{itemize}
    
Por fim, o autor reitera que a utilização destas ferramentas de Inteligência Artificial possuiu caráter consultivo e de apoio estrutural. Dessa forma, o autor atesta ter revisado integralmente o material gerado e assume a responsabilidade total, irrestrita e exclusiva sobre todo o conteúdo intelectual, analítico e técnico apresentado neste documento, bem como pela originalidade da pesquisa e pela garantia de ausência de plágio.